\section{Symmetry and special slices}\label{sec:symmetry}

This section describes a discrete symmetry of KRMHD and two slices with special transport properties. The symmetry is the reflection of the plasma in a plane perpendicular to the guide field. It exchanges the slices $u$ and $-u$, divides the invariants and diagnostics of the previous sections into even and odd ones, and gives exact statements about what the injection of $\Gamma$ can change (sections~\ref{sym:map}--\ref{sym:parity}). The special slices are those with $su=\pm1$, each of which is advected by a single Elsasser field (section~\ref{sym:vA}). Section~\ref{sym:fluid} relates them, and the slice $u=0$, to the slow waves and the entropy mode of the fluid theory. Throughout this section $\langle\cdot\rangle$ denotes the average over the periodic box, $\mathrm E$ an ensemble mean, and $S$ the reflection of the velocity variable, $(SG)(u)=G(-u)$. In Hermite variables $S=\mathrm{diag}\big((-1)^m\big)$, because $H_m(-u)=(-1)^mH_m(u)$.

\subsection{The mirror map}\label{sym:map}

We define the mirror map $\mathcal M$ by
\begin{equation}\label{eq:mirror}
  \Phi'(\boldsymbol r)=\Phi(\tilde{\boldsymbol r}),\qquad
  \Psi'(\boldsymbol r)=-\Psi(\tilde{\boldsymbol r}),\qquad
  G'(\boldsymbol r,u)=G(\tilde{\boldsymbol r},-u),\qquad
  \tilde{\boldsymbol r}=(x,y,-z),
\end{equation}
where primes denote the mirrored state and the last relation applies to both compressive branches and to the sector $\Lambda=\infty$. In Elsasser and Hermite variables, $z^{\pm\prime}(\boldsymbol r)=z^\mp(\tilde{\boldsymbol r})$ and $g'_m(\boldsymbol r)=(-1)^mg_m(\tilde{\boldsymbol r})$, that is, $g'=S\tilde g$, where a tilde on a field denotes $\tilde f(\boldsymbol r)=f(\tilde{\boldsymbol r})$. The map is the reflection of the plasma in the plane $z=0$. Under this reflection the flow $\uE$ transforms as a polar vector and the magnetic perturbation $\boldsymbol b_\perp$ as an axial vector, the guide field is unchanged, and $v_\parallel$ changes sign. The map is an involution and reverses the direction of propagation along the guide field: the mirror image of a $z^+$ wave packet, which propagates toward $-z$, is a $z^-$ wave packet, which propagates toward $+z$.

The map preserves the equations of KRMHD. Let $\gradpar'=\partial_z+\{\Psi',\cdot\}$ be the derivative along the mirrored field lines. For every field $f$,
\begin{equation}\label{sym:gradpar}
  \{\Phi',\tilde f\}=\widetilde{\{\Phi,f\}},\qquad
  \gradpar'\tilde f=-\widetilde{\gradpar f},
\end{equation}
because the bracket involves only $x$ and $y$, while $\partial_z$ and $\Psi$ both change sign. Every term of the induction equation in (\ref{eq:rmhd}) therefore changes sign and every term of the vorticity equation keeps its sign, so both equations are invariant. In the kinetic equation (\ref{eq:branch}) the streaming term changes sign twice, through $\gradpar$ and through $L$: since $F_M$ is even and $M_0[SG]=M_0[G]$, the operator (\ref{eq:L}) satisfies
\begin{equation}\label{sym:SLS}
  SLS=-L .
\end{equation}
In Hermite variables this is $SAS=-A$, which holds for the zero closure and for the copy closure, because in both the matrix $A$ has non-zero entries only next to the diagonal. The equations therefore commute with $\mathcal M$ when the collision operator commutes with $S$ and every dissipation term is a Fourier multiplier that depends on $k_z$ only through $k_z^2$, applied to $\Phi$ and to $\Psi$ without mixing them, and to $g$ through a matrix in $m$ that commutes with $S$, for example the same multiplier for every $g_m$. Operators that are diagonal in $m$, such as GANDALF's hypercollisions $\nu(m/M)^n$, commute with $S$. GANDALF's perpendicular hyper-dissipation and its parallel dissipation $\eta_z$ act in the same way on $z^+$, $z^-$ and every $g_m$, with multipliers that depend on $k_\perp$ and on $k_z^2$, and so satisfy the condition. We checked that one GANDALF time step commutes with $\mathcal M$ to $6\times10^{-16}$, for both time integrators and both closures, at $\Lambda=\pm\sqrt5$, with hypercollisions, hyper-resistivity and $\eta_z$ switched on (section~\ref{sec:numerics}).

For $\Lambda\notin[0,1]$ the map acts on the slices as follows. By (\ref{sym:SLS}), $SF_u$ is a Case--Van Kampen mode with phase velocity $-u$, and $M_0[SF_u]=M_0[F_u]=1$, so $SF_u=F_{-u}$. The reflection preserves the free-energy product, $\fe{Sf}{Sh}=\fe{f}{h}$, and $\varrho$ is even, because $\eps(-u+\ii0)$ is the complex conjugate of $\eps(u+\ii0)$. Hence $\Tr[SG](u)=\Tr[G](-u)$, and the unit-weight amplitudes (\ref{eq:aunit}) of the mirrored state are
\begin{equation}\label{sym:slices}
  a_u[\mathcal MX](\boldsymbol r)=a_{-u}[X](\tilde{\boldsymbol r}),
\end{equation}
where $X=(\Phi,\Psi,G)$ denotes the state. The mirror exchanges slice $u$ with slice $-u$, and with them their advecting potentials: by (\ref{eq:mirror}), $\Phi'+su\Psi'$ is the reflection of $\Phi-su\Psi$, the potential of slice $-u$ in (\ref{eq:slice}). The slice energies $\langle a_u^2\rangle$ and $\langle a_{-u}^2\rangle$ are exchanged, and the slice $u=0$, which is advected by $\uE$ alone, is mapped to itself. In a truncated hierarchy the slice speeds come in pairs $\pm u_j$, $S$ maps the eigenvector of $u_j$ to that of $-u_j$, and for even $M$ the zero-speed slice is mapped to itself (section~\ref{sec:truncated}).

A statistic $O$, that is, a functional of the state at one time or of its history, is mirror-even if $O(\mathcal MX)=O(X)$ for every $X$ and mirror-odd if $O(\mathcal MX)=-O(X)$. By (\ref{sym:gradpar}), the quadratic single-point statistic $\langle g_m\gradpar^qg_n\rangle$ is multiplied by $(-1)^{q+m+n}$ under $\mathcal M$. By (\ref{sym:SLS}) and (\ref{sym:slices}), the invariant $\Qphi$ of (\ref{eq:Qphi}) is mapped to $Q_{\check\varphi}$ with $\check\varphi(u)=\varphi(-u)$, and the Casimir $\int\mathcal F(a_u)\,\dd^3\boldsymbol r$ of slice $u$ to that of slice $-u$. This gives the classification below. In it, $\Wc(m)=p_m\langle g_m^2\rangle$ is the Hermite spectrum of the free energy; $\Gamma(m)=2p_mA_{m,m+1}\langle g_mg_{m+1}\rangle$ is the contribution of the pair $(m,m+1)$ to $\langle\fe{G}{LG}\rangle=\sum_m\Gamma(m)$ (section~\ref{sec:decorrelation}), which is the Hermite form (B32) of CMM's $\Gamma$ up to normalisation; $\Pi_{m+1/2}$ is the field-aligned Hermite flux (\ref{eq:fluxCD}); $\eps_\nu$ is the collisional dissipation rate; and $E^\pm=\langle|\nabla_\perp z^\pm|^2\rangle/4$ are the energies of the Elsasser fields.
\begin{itemize}
  \item Mirror-even: $\Wc(m)$ and $\Wc$; $\Pi_{m+1/2}$; the $k_z$-based Hermite flux summed separately over the Fourier modes in which it is positive and over those in which it is negative (the mirror permutes the per-mode fluxes, $\boldsymbol k\to\tilde{\boldsymbol k}$); the energies of the phase-mixing and phase-unmixing parts in the $\mathrm{sgn}\,k_z$ decomposition (section~\ref{sec:decorrelation}); $\eps_\nu$; the spectrum $\Wc(k_\perp)$; $Q_{u^{2n}}$; the energy $\langle a_0^2\rangle$ of the slice $u=0$, which gives the imaginary part of the complex helicity (\ref{eq:Hcomplex}); the integral of $\langle a_u^4\rangle$ over all slices (the sum over the slices in a truncated hierarchy); $\langle g_m^4\rangle$; $E^++E^-$.
  \item Mirror-odd: $\Gamma\propto Q_u$ and every $\Gamma(m)$; $\Hph\propto Q_{1/u}$; $Q_{u^3}$; $\sum_{\boldsymbol k}\mathrm{sgn}(k_z)\,|g_m(\boldsymbol k)|^2$; the cross-helicity $E^+-E^-$.
\end{itemize}
Every state that is its own mirror image therefore has $\Gamma=\Hph=0$ and $\Gamma(m)=0$ for every $m$, whatever its Hermite spectrum; proposition~\ref{prop:nofjortoft} uses this. That $\Hph$ is odd under $v_\parallel\to-v_\parallel$ is also explicit in CMM's eq.~(2.19), which writes the real part of the gyrokinetic helicity as an integral of the product of the parts of their function $\psi_s$ (their eq.~2.18) that are even and odd in $v_\parallel$.

\subsection{Mirror-symmetric forced ensembles}\label{sym:ensembles}

The density $\fe{G}{LG}$ of $\Gamma$ contains $g_0$ only in the term $\sqrt2\,p_0g_0g_1$, so a correlation between the forcings of $g_0$ and $g_1$ injects $\Gamma$ (section~\ref{sec:slices}). We consider the pair forcing
\begin{equation}\label{sym:pair}
  F_\phi[f_0]=\big(f_0,\;r\,\mathcal R_\phi f_0,\;0,\;0,\;\dots\big),\qquad
  \mathcal R_\phi=\cos\phi+\sin\phi\,\mathcal J_x ,
\end{equation}
added to the right-hand side of the Hermite hierarchy of one compressive branch. Here $f_0(\boldsymbol r,t)$ is a real random field, $r>0$ and $\phi$ are constants, and $\mathcal J_x$ is the real antisymmetric operator with Fourier multiplier $\ii\,\mathrm{sgn}\,k_x$ (zero at $k_x=0$), which is the transform $\mathsf H$ of section~\ref{sec:slices} acting in $x$. Because $\langle f_0\,\mathcal J_xf_0\rangle=0$, the correlation $\langle F_0F_1\rangle=r\cos\phi\,\langle f_0^2\rangle$ comes from the $\cos\phi$ part alone, and only this part injects $\Gamma$. Under the mirror, both parts of the forcing of $g_1$, the $\cos\phi$ part and the $\sin\phi$ part, change sign relative to the forcing of $g_0$. The mirror reverses the sign of $g_1$ and maps $f_0$ to $\tilde f_0$, and $\mathcal J_x$ commutes with the reflection $z\to-z$ because its multiplier does not depend on $k_z$. Hence
\begin{equation}\label{sym:MF}
  \mathcal MF_\phi[f_0]=F_{\phi+\pi}[\tilde f_0] .
\end{equation}
The mirror image of the pair forcing at phase $\phi$ is the pair forcing at phase $\phi+\pi$, which is the same forcing with $r$ replaced by $-r$. If the phase is implemented by multiplying by $e^{\ii\phi}$ the stored coefficients of a real-to-complex transform, which holds the modes with $k_x\ge0$, the real field obtained by Hermitian completion is $r\mathcal R_\phi f_0$ at every $k_x\neq0$, so its $\sin\phi$ part is a Hilbert transform in $x$, not a transform in $z$. On the plane $k_x=0$ the multiplied coefficients do not satisfy the reality condition unless $\sin\phi=0$; the analogue of (\ref{sym:MF}) for the stored coefficients, with $e^{\ii\phi}$ in place of $\mathcal R_\phi$, holds in either case. GANDALF's \texttt{force\_hermite\_\allowbreak moments} with \texttt{forced\_moments=(0,1)} adds the same pattern to $g_0$ and $g_1$, which is (\ref{sym:pair}) with $r=1$ and $\phi=0$.

\begin{theorem}[Mirror symmetry of forced ensembles]\label{thm:mirror}
Consider the Alfv\'enic fields and one compressive branch (the other compressive sectors are passive and do not affect them), in the continuum or in a Fourier--Hermite truncation with either closure, with collisions and dissipation that satisfy the conditions of section~\ref{sym:map}; in a time-discrete system, let the time step commute with $\mathcal M$. Let $(z^+,z^-)$ be forced by a random forcing $F_A=(F^+,F^-)$ and the branch by the pair forcing (\ref{sym:pair}), and assume that the solution is a deterministic, measurable function of the initial state and of the realisation of the forcing. Write $X^\phi$ for the solution, a random process in time, and $\mathrm E_\phi$ for the mean over its law. Assume that
\begin{enumerate}
  \item[(a)] the law of $F_A$ is invariant under $\mathcal M$, that is, $(\tilde F^-,\tilde F^+)$ has the law of $(F^+,F^-)$;
  \item[(b)] the law of $f_0$ is invariant under $z\to-z$, that is, $\tilde f_0$ has the law of $f_0$;
  \item[(c)] the initial state, $F_A$ and $f_0$ are independent, and the initial laws at the phases $\phi$ and $\phi+\pi$ are mirror images, $\mathrm{law}\,X^{\phi+\pi}(0)=\mathrm{law}\,\mathcal MX^\phi(0)$.
\end{enumerate}
Then $\mathrm{law}\,X^{\phi+\pi}=\mathrm{law}\,\mathcal MX^\phi$, and for every statistic $O$ whose expectation exists, of any order,
\begin{equation}\label{sym:Ephi}
  \mathrm E_{\phi+\pi}[O]=\mathrm E_\phi[O]\quad\text{if $O$ is mirror-even},\qquad
  \mathrm E_{\phi+\pi}[O]=-\mathrm E_\phi[O]\quad\text{if $O$ is mirror-odd}.
\end{equation}
\end{theorem}

\begin{proof}
Write the solution as $X^\phi=\mathcal Z_\phi\big(X(0),F_A,f_0\big)$. The forcing of $g$ transforms as $g$ does, so if $X$ solves the equations with initial state $X(0)$ and forcings $F_A$ and $F_\phi[f_0]$, then $\mathcal MX$ solves them with initial state $\mathcal MX(0)$ and forcings $\mathcal MF_A=(\tilde F^-,\tilde F^+)$ and $\mathcal MF_\phi[f_0]$, because every term of the equations commutes with $\mathcal M$ (section~\ref{sym:map}). By (\ref{sym:MF}) and the uniqueness of the solution,
\begin{equation*}
  \mathcal Z_{\phi+\pi}\big(\mathcal MX(0),\mathcal MF_A,\tilde f_0\big)=\mathcal M\,\mathcal Z_\phi\big(X(0),F_A,f_0\big).
\end{equation*}
By (a)--(c), the inputs $\big(X^{\phi+\pi}(0),\allowbreak F_A,\allowbreak f_0\big)$ of the run at $\phi+\pi$ have the joint law of $\big(\mathcal MX^\phi(0),\allowbreak\mathcal MF_A,\allowbreak\tilde f_0\big)$, so $X^{\phi+\pi}$ has the law of $\mathcal MX^\phi$. Therefore $\mathrm E_{\phi+\pi}[O(X)]=\mathrm E_\phi[O(\mathcal MX)]$, which is (\ref{sym:Ephi}).
\end{proof}

The GANDALF routines \texttt{force\_alfven\_\allowbreak modes} and \texttt{force\_alfven\_modes\_\allowbreak balanced} satisfy (a). The first adds the same noise to $z^+$ and $z^-$, so that only $\Phi$ is forced, and the second adds identically distributed noise to $z^+$ and $z^-$ (independent by default). In both routines the set of forced modes depends on $k_z$ only through $|\boldsymbol k|$ or $|k_z|$, and the rule that makes the coefficients on the plane $k_x=0$ real is applied index by index, so the law of the noise is invariant under $k_z\to-k_z$. A routine that forces a prescribed list of modes, such as \texttt{force\_alfven\_modes\_\allowbreak specific}, satisfies (a) only if the set of forced wavevectors $\pm\boldsymbol k$ is symmetric under $k_z\to-k_z$; the six modes that the documentation of that routine gives for the stirring of the original GANDALF code are not. The noise used by \texttt{force\_hermite\_\allowbreak moments} is generated in the same way and satisfies (b). Hypothesis (c) holds when both runs start from the same mirror-invariant law, for example from zero fields. In that case the law of the Alfv\'enic fields is itself invariant under $\mathcal M$, and every mirror-odd statistic of the Alfv\'enic fields alone, such as the cross-helicity $E^+-E^-$, has zero mean; in this sense such ensembles are balanced. When the two runs start from the same state and that state is not mirror-symmetric, (c) fails, and the theorem applies only to the extent that the statistics have lost the memory of the initial state. The proof uses only the equivariance of the dynamics, and we checked it pathwise: in a GANDALF integration with pair forcing, hyper-resistivity and hypercollisions, a run started from $\mathcal MX$ with mirrored noise and phase $\phi+\pi$ remained the mirror image of the run at $\phi$ to $3.7\times10^{-15}$ after 200 steps of the IMEX scheme, with equal $\Wc(m)$ and opposite $\Gamma$ (section~\ref{sec:numerics}). We did not compute ensemble averages; the theorem rests on the equivariance, which these checks test directly.

The theorem has a direct consequence for the injection of $\Gamma$. The phases $\phi=0$ and $\phi=\pi$ correspond to forcing $g_1$ with $+rf_0$ and with $-rf_0$: the two ensembles receive $\Gamma$ of opposite signs and the same free energy. By (\ref{sym:Ephi}) they have the same expectation of every mirror-even statistic, at every order and with all dissipation and collision terms included. When the two runs start from the same mirror-invariant initial law, the sign of the $\Gamma$ injection therefore cannot change the mean of $\Wc(m)$, of the field-aligned or $k_z$-based Hermite fluxes, of $\eps_\nu$ or of any other mirror-even statistic. Only mirror-odd statistics, beginning with $\Gamma$ itself, distinguish $\phi$ from $\phi+\pi$. If, in addition, $f_0$ has zero mean, more follows for statistics that are quadratic in $g$, such as $\Wc(m)$, $\eps_\nu$ and $\Pi_{m+1/2}$. For given Alfv\'enic fields $g$ is linear in its forcing, and $f_0$ is independent of the Alfv\'enic forcing, so the part of the mean of such a statistic that comes from the correlation between the forcings of $g_0$ and $g_1$ is linear in $r$, and (\ref{sym:Ephi}) makes it vanish for mirror-even statistics. Their means then equal those obtained when $g_0$ and $g_1$ are forced by independent patterns, each with the covariance it has in the pair forcing.

A correlation through an operator that is odd under $z\to-z$ behaves differently. If $g_1$ is forced with $r\mathcal J_zf_0$, where $\mathcal J_z$ has the multiplier $\ii\,\mathrm{sgn}\,k_z$, the mirror maps the forcing to $(\tilde f_0,\,r\mathcal J_z\tilde f_0,\,0,\dots)$, so this kind of correlation is mirror-even. It injects no $\Gamma$, since $\langle f_0\,\mathcal J_zf_0\rangle=0$, but it correlates $g_0$ with $\partial_zg_1$, so the injected pattern carries a Hermite flux $\Pi_{1/2}$ whose sign changes with that of $r$. The forcings with $r$ and $-r$ are then not mirror images of each other, and the theorem does not constrain the difference between them.

\subsection{A parity lemma under frozen-line dynamics}\label{sym:parity}

Theorem~\ref{thm:mirror} concerns ensembles that are symmetric under $\mathcal M$. Under the hypotheses of the frozen-line theorem a statement of the same kind holds for each Alfv\'enic realisation separately, for quadratic single-point statistics. Let $e_m$ be the Hermite basis functions of section~\ref{sec:slices}, so that $g_n(e_m)=\delta_{nm}$. An initial condition $g(0)=f\,(e_0+re_1)$, that is, $g_0=f$, $g_1=rf$ and $g_m=0$ for $m\ge2$, with $f$ a real field and $r$ real, is the $\Gamma$-injecting ($\cos\phi$) part of a pair correlation. It carries $Q_u=\sqrt2\,p_0r\int f^2\,\dd^3\boldsymbol r$, which changes sign with $r$, while $Q_1=(p_0+r^2)\int f^2\,\dd^3\boldsymbol r$, which is proportional to its free energy, does not depend on the sign of $r$.

\begin{lemma}[Parity lemma]\label{lem:parity}
Assume (\ref{eq:hyp}), in the continuum or in a truncated hierarchy with either closure, with a collision operator $\mathcal C$ whose matrix in the Hermite basis is real and commutes with $S$ (for example, an operator that is diagonal in $m$). Let $g(0)=f\,(e_0+re_1)$, where $f$ is a smooth real field and $r$ is real. Then for $0\le t\le T$ every mirror-even quadratic single-point statistic of $g(t)$, that is, every linear combination of the $\langle g_m\gradpar(t)^qg_n\rangle$ with $q+m+n$ even, is an even function of $r$. This includes $\Wc(m,t)$, $\eps_\nu(t)$ and the field-aligned flux $\Pi_{m+1/2}(t)$. The mirror-odd statistics, such as $\Gamma(m,t)$, have in general a part linear in $r$.
\end{lemma}

\begin{proof}
By corollary~\ref{cor:singlepoint}, each such statistic equals the corresponding statistic of $h(t)=e^{-tL_0}g(0)$, computed with $\gradpar(0)$, so it suffices to treat the frozen linear problem. The operator $K_\parallel(0)=-\ii\gradpar(0)$ of corollary~\ref{cor:kpar} is self-adjoint and anticommutes with complex conjugation, because $\gradpar(0)$ is real. The spectral measure $\mu_f(I)=\langle f\,E_0(I)f\rangle$ of the real field $f$, where $E_0$ is the spectral measure of $K_\parallel(0)$ as in section~\ref{sec:lagrangian}, is therefore symmetric, $\mu_f(-I)=\mu_f(I)$. In the spectral representation of $K_\parallel(0)$, $\gradpar(0)$ is multiplication by $\ii\kappa$, and $e^{-tL_0}$ acts on the Hermite coefficients as the matrix $\Theta(\kappa)=\exp[-t(\ii s\kappa A+\mathcal C)]$, where $A$ is the ladder matrix and $\mathcal C$ now denotes the matrix of the collision operator (the propagator replaces the exponential if $\mathcal C$ depends on $t$). Hence
\begin{equation*}
  \langle h_m\gradpar(0)^qh_n\rangle=\int\mu_f(\dd\kappa)\,(\ii\kappa)^q\,
  \overline{\big[\Theta_{m0}+r\Theta_{m1}\big]}\,\big[\Theta_{n0}+r\Theta_{n1}\big],
\end{equation*}
and the part linear in $r$ has the integrand $(\ii\kappa)^q\big[\overline{\Theta_{m0}}\,\Theta_{n1}+\overline{\Theta_{m1}}\,\Theta_{n0}\big]$. Let $D=\mathrm{diag}\big((-\ii)^m\big)$. Because $A$ has non-zero entries only next to the diagonal, $\hat A=D^{-1}(\ii A)D$ is real, with $\hat A_{m,m+1}=A_{m,m+1}$ and $\hat A_{m+1,m}=-A_{m+1,m}$, and $S\hat AS=-\hat A$. Because $\mathcal C$ is real and commutes with $S$, its entries vanish unless $m+n$ is even, so $\hat{\mathcal C}=D^{-1}\mathcal CD$ is real and commutes with $S$. Hence
\begin{equation}\label{sym:similarity}
  \Theta(\kappa)=D\,\Xi(\kappa)\,D^{-1},\qquad
  \Xi(\kappa)=\exp[-t(s\kappa\hat A+\hat{\mathcal C})]\ \ \text{real},\qquad
  \Xi(-\kappa)=S\,\Xi(\kappa)\,S,
\end{equation}
so that $\Theta_{ab}=\ii^{\,b-a}\,\Xi_{ab}$ and $\Xi_{ab}(-\kappa)=(-1)^{a+b}\,\Xi_{ab}(\kappa)$. Then $\overline{\Theta_{m0}}\,\Theta_{n1}=\ii^{\,m-n+1}\,\Xi_{m0}\Xi_{n1}$ and $\overline{\Theta_{m1}}\,\Theta_{n0}=\ii^{\,m-n-1}\,\Xi_{m1}\Xi_{n0}$, and the integrand acquires the factor $(-1)^{q+m+n+1}$ under $\kappa\to-\kappa$. When $q+m+n$ is even, which is the condition for $\langle g_m\gradpar^qg_n\rangle$ to be mirror-even, the integrand is odd, and its integral against the symmetric measure $\mu_f$ vanishes.
\end{proof}

The lemma holds for every Alfv\'enic realisation that satisfies (\ref{eq:hyp}); no symmetry of the Alfv\'enic fields is assumed. The similarity (\ref{sym:similarity}) is where the structure of the hierarchy enters: the streaming couples only neighbouring Hermite moments, and the collisions do not couple moments of opposite parity.

We tested the lemma with GANDALF's right-hand side at $M=8$, $\Lambda=\sqrt5$, $\beta_i=1$ and $r=0.8$, in a random ideal Alfv\'enic flow without mirror symmetry. For a set of numbers $X_m$ computed in the runs with $+r$ and $-r$, we use the relative difference $\sum_m|X_m^+-X_m^-|/\sum_m|X_m^++X_m^-|$. In the frozen problem, whose discrete $\gradpar(0)$ is a real antisymmetric matrix on the dealiased fields, this difference is at most $3.4\times10^{-15}$ at every resolution, both for $X_m=\langle g_m^2\rangle=\Wc(m)/p_m$ and for $X_m=\Pi_{m+1/2}$. In the nonlinear integrations, advanced with RK4, the difference for $\langle g_m^2\rangle$ at $t=0.062$ is $2.8\times10^{-5}$, $2.6\times10^{-6}$, $1.1\times10^{-6}$ and $2.0\times10^{-7}$ on $16^3$, $24^3$, $32^3$ and $48^3$ grids, and the difference for $\Pi_{m+1/2}$ falls from $3.2\times10^{-4}$ on the $16^3$ grid to $4.7\times10^{-6}$ on the $48^3$ grid, while $\Gamma$ takes opposite values in the two runs. Two variants that the lemma does not cover behave differently. With the mirror-even correlation $g_1=\pm r\mathcal J_zf$ of section~\ref{sym:ensembles} in place of $g_1=\pm rf$, the difference for $\langle g_m^2\rangle$ is $0.37$ at $t=0.062$ and $0.43$ at $t=0.124$ on $16^3$, $24^3$ and $32^3$ grids. With a perpendicular damping term $-4\times10^{-6}k_\perp^4g$, which violates (H2), it is $2.7\times10^{-2}$, $3.5\times10^{-2}$ and $3.5\times10^{-2}$ at $t=0.124$ on $16^3$, $24^3$ and $32^3$ grids, and does not decrease with resolution. Whether such a departure persists as the dissipation coefficient tends to zero is not addressed here (section~\ref{sec:discussion}).

Together with corollary~\ref{cor:forcing}, the lemma extends to forcing. Let $g$ be forced by $F=f\,(e_0+re_1)$, where $f$ is a real random field with zero mean that is white in time and independent of $g(0)$ and of the Alfv\'enic history, and $r$ is real. For every realisation of the Alfv\'enic fields that satisfies (H1) and (H3) of (\ref{eq:hyp}), with $g$ satisfying (H2) apart from the source, corollary~\ref{cor:forcing} writes the mean over the forcing of $g$ of a quadratic single-point statistic as the sum of the contribution of $g(0)$, which does not depend on $r$, and an integral over injection times. The integrand is the contribution of the injection at that time, computed in the frozen-line problem with the field lines frozen at the injection time. The lemma, with $\gradpar(t')$ in place of $\gradpar(0)$, applies to each injection. The mean over the forcing of $g$ of every mirror-even quadratic single-point statistic is therefore even in $r$. It also equals the mean obtained when $g_0$ and $g_1$ are forced by independent patterns, each with the covariance it has in the pair forcing, because for each injection the part linear in $r$ vanishes and the remaining parts are those of separate injections into $g_0$ and $g_1$. For forcing of $g$ that is white in time and independent of the Alfv\'enic history, the correlation that injects $\Gamma$ therefore has no effect on the mean, over the forcing of $g$, of any mirror-even quadratic single-point statistic, for every Alfv\'enic realisation that satisfies (H1) and (H3). Forcing with a finite correlation time is not covered: the mean response to two correlated injections at different times contains the relabelling by the flow between them (corollary~\ref{cor:forcing}), to which the lemma does not apply, and it need not be even in $r$.

The lemma and its extension concern quadratic statistics. At fourth order a mirror-even statistic can distinguish a correlated injection from an uncorrelated one with the same variances. By (\ref{eq:sliceforcing}), an injection of $g_0=f$ and $g_1=rf$ loads slice $u$ with $a_u=\varrho(u)^{1/2}(p_0+\delta_u)f$, where $\delta_u=\sqrt2\,ru$, while an injection of $g_0=f$ and $g_1=rf'$, with $f'$ independent of $f$ and identically distributed, loads it with $a_u=\varrho(u)^{1/2}(p_0f+\delta_uf')$. The two injections give the same means of the mirror-even second-order single-point statistics, such as $\Wc(m)$ and the free energy, which is proportional to $p_0+r^2$ per unit variance of $f$. Their mean slice energies differ by $2\varrho(u)p_0\delta_u$ per unit variance, which is odd in $u$; its first moment in $u$, $\sqrt2\,p_0r$, is the $Q_u$ of the correlated state. For statistically homogeneous Gaussian $f$ and $f'$ of unit variance, the mirror-even fourth-order statistic of the slice pair $\pm u$ differs by
\begin{equation}\label{sym:fourth}
  \mathrm E\big\langle a_u^4+a_{-u}^4\big\rangle_{\rm corr}-\mathrm E\big\langle a_u^4+a_{-u}^4\big\rangle_{\rm uncorr}
  =24\,\varrho(u)^2\,p_0^2\,\delta_u^2 .
\end{equation}
In a truncated hierarchy the same algebra holds node by node, with $\varrho(u)$ replaced by the Christoffel number $\lambda_j$ of the node (zero closure; section~\ref{sec:truncated}); at $\Lambda=\sqrt5$ and $r=1$ the sum over the slices of $\mathrm E\,a_j^4$ is $55\%$ larger for the correlated state at $M=8$ and $61\%$ larger at $M=64$. The difference (\ref{sym:fourth}) is even in $r$, so fourth-order slice statistics detect the presence of the correlation that injects $\Gamma$ but not its sign. In a mirror-symmetric ensemble no mirror-even statistic, of any order, separates $\phi=0$ from $\phi=\pi$ in the mean (theorem~\ref{thm:mirror}), whereas the mirror-odd statistics, beginning with $\Gamma$, separate them at second order.

\subsection{The slices at \texorpdfstring{$v_\parallel=\pm v_A$}{v = +-vA}}\label{sym:vA}

\begin{proposition}\label{prop:elsasser}
Under the hypotheses of theorem~\ref{thm:slice}, the slices with $su=-1$ and $su=+1$, that is, with $v_\parallel=u\vth=-v_A$ and $+v_A$, obey
\begin{equation}\label{sym:elsasser}
  (\partial_t-\partial_z)\,a_u+\{z^-,a_u\}=0\quad(su=-1),\qquad
  (\partial_t+\partial_z)\,a_u+\{z^+,a_u\}=0\quad(su=+1).
\end{equation}
Each is advected by a single Elsasser field and streams with the other: the slice $su=-1$ moves toward $-z$ with $z^+$ and is advected by $z^-$ alone, and the slice $su=+1$ moves toward $+z$ with $z^-$ and is advected by $z^+$ alone. In a truncated hierarchy the same holds for a slice of theorem~\ref{thm:trunc} whose speed $su_j$ equals $-1$ or $+1$.
\end{proposition}

\begin{proof}
By (\ref{eq:elsasserpot}), the advecting potential $\Phi+su\Psi=\frac12\big[(1+su)z^++(1-su)z^-\big]$ is $z^-$ at $su=-1$ and $z^+$ at $su=+1$. Substituting in (\ref{eq:slice}) gives (\ref{sym:elsasser}).
\end{proof}

The slice with $su=-1$ is a passive counterpart of the Elsasser field $z^+$. By (\ref{eq:rmhd}), the Elsasser vorticity $\nabla_\perp^2z^+$ obeys
\begin{equation}\label{sym:vorticity}
  \big(\partial_t-\partial_z+\{z^-,\cdot\}\big)\,\nabla_\perp^2z^+
  =\{\partial_xz^+,\partial_xz^-\}+\{\partial_yz^+,\partial_yz^-\},
\end{equation}
and the operator on the left is the transport operator of (\ref{sym:elsasser}) at $su=-1$. The right-hand side of (\ref{sym:vorticity}) arises because $\boldsymbol Z^\pm=\zhat\times\nabla_\perp z^\pm$ are vector fields. It equals $(\boldsymbol Z^-\cdot\nabla)\nabla_\perp^2z^+-\zhat\cdot\nabla\times\big[(\boldsymbol Z^-\cdot\nabla)\boldsymbol Z^+\big]$, the difference between the advection of the vorticity of $\boldsymbol Z^+$ and the curl of the advection term in the Elsasser equation $\partial_t\boldsymbol Z^+-\partial_z\boldsymbol Z^++(\boldsymbol Z^-\cdot\nabla)\boldsymbol Z^+=-\nabla p$, and it vanishes when either Elsasser field vanishes. The slice is transported by the same operator without this term, and the slice with $su=+1$ is in the same way a passive counterpart of $z^-$. Like the Elsasser fields, each of the two slices is distorted only by the counter-propagating field. When $z^-=0$, any $z^+=f(x,y,z+t)$ is an exact solution of (\ref{eq:rmhd}). The slice with $su=-1$ then obeys $(\partial_t-\partial_z)a_u=0$ and translates rigidly with $z^+$, while every other slice is advected by $\frac12(1+su)z^+$, streams at a different speed, and is in general distorted.

In the continuum the two slices exist at every $\beta_i$, at $u=\pm\beta_i^{-1/2}$. They lie in the bulk of the distribution for $\beta_i\gtrsim1$ and in its tail for $\beta_i\ll1$. In a truncated hierarchy the slice speeds are the eigenvalues $u_j$ of $A$ (theorem~\ref{thm:trunc}), so exact slices with $su=\pm1$ exist only when $\beta_i=u_j^{-2}$ for some $j$. We checked proposition~\ref{prop:elsasser} against GANDALF with $M=8$, $\Lambda=\sqrt5$ and $\beta_i=2.486$, for which $su_j=-1$ at the eigenvalue $u_j=-0.634$. GANDALF takes $\Lambda$ and $\beta_i$ as separate inputs (section~\ref{sec:model}). This pair is not a KRMHD branch, because at this $\beta_i$ the value $\Lambda^+$ of (\ref{eq:Lambdapm}) lies between $1.40$ and $1.48$ for every $\tau/Z$ and $\Lambda^-$ is negative. The test does not depend on this, because the proposition holds for every $\Lambda\notin[0,1]$ and every $s$. The projection of GANDALF's right-hand side onto this slice equals $\partial_za_j-\{z^-,a_j\}$ to $5.5\times10^{-16}$ and does not change when $z^+$ is replaced by an unrelated field. At $M=9$, $\Lambda=-\sqrt5$ and $\beta_i=0.830$, which is the branch $\Lambda^-$ at $\tau/Z=1.06$, the slice with $u_j=-1.098$ has $su_j=-1$, and the same holds to $1.6\times10^{-15}$. A GANDALF integration with the Lawson-RK4 scheme at the first set of parameters, from $z^-=0$ to $t=0.5$, kept $z^-=0$, translated $z^+$ rigidly, and translated this slice rigidly to $1.4\times10^{-14}$. Every other slice departed from rigid streaming by 1.35 to 1.4 in relative norm, close to the value $\sqrt2$ for a field uncorrelated with its rigidly streamed initial state.

In the fluid limit the corresponding property is known at high $\beta$. S09 (\S2.5) showed from their eq.~(30) that only for $\beta\gg1$ do the slow waves interact exclusively with the counter-propagating Alfv\'en waves, so that setting one Elsasser potential to zero gives finite-amplitude slow-wave packets; their footnote~15 adds that such solutions exist in the Lagrangian frame of the Alfv\'en waves. \citet[\S2.4]{LithwickGoldreich2001} described the slow waves of incompressible MHD, which is the limit $\beta\gg1$, as travelling along the field lines at the Alfv\'en speed and being distorted by interactions with oppositely directed Alfv\'en waves. At finite $\beta$ the fluid slow waves are also scattered by the co-propagating Alfv\'en waves (S09, \S\S2.5 and~2.7). What is new here is that the collisionless compressive sector contains, at every $\beta_i$, two kinetic slices, at $v_\parallel=\pm v_A$, that are each advected by a single Elsasser field, a property that the fluid slow waves have only for $\beta\gg1$.

\subsection{Fluid slow waves and the entropy mode}\label{sym:fluid}

In the collisional limit the compressive fluctuations obey fluid equations, and these have the form of the slice equation. S09 wrote the slow-wave equations, their eqs~(26)--(27), as equations for the generalised Elsasser fields of their eq.~(29), which in units $v_A=1$ are $z_\parallel^\pm=u_\parallel\pm(1+v_A^2/c_s^2)^{1/2}\,\delta B_\parallel/B_0$, with $c_s$ the sound speed. With $c=(1+v_A^2/c_s^2)^{-1/2}$, the phase speed of the slow waves in units of $v_A$, these equations (their eq.~30) and the equation $\dd\delta s/\dd t=0$ for the entropy mode $\delta s$ (their eq.~23) read
\begin{equation}\label{eq:fluidslices}
  \partial_tz_\parallel^\pm+\{\Phi\mp c\Psi,\,z_\parallel^\pm\}\mp c\,\partial_zz_\parallel^\pm=0,\qquad
  \partial_t\,\delta s+\{\Phi,\delta s\}=0 .
\end{equation}
Equations (\ref{eq:fluidslices}) are the slice equation (\ref{eq:slice}) at $su=\mp c$ and at $u=0$. This is a formal correspondence between equations: the slow-wave fields and the entropy mode obey the equation of a slice with the same speed. It is not a derived collisional limit of the kinetic slices. S09 showed that KRMHD reduces to these fluid equations when collisions dominate (their \S6.1 and appendix~D), but in the collisional system the continuum of slice speeds is replaced by the three characteristic speeds $-c$, $0$ and $c$, and we have not shown how the kinetic slices pass over into $z_\parallel^\pm$ and $\delta s$. For $\beta\gg1$, $c\to1$ and the fluid fields move to $su=\mp1$, the speeds of the slices of proposition~\ref{prop:elsasser}; this is the high-$\beta$ statement of S09 and of \citet{LithwickGoldreich2001} recalled in section~\ref{sym:vA}. The slice $u=0$ is advected by $\uE$ alone at every $\beta_i$, as the entropy mode is. At low $\beta_i$, where $s\ll1$, every slice with $u$ of order unity has $|su|\ll1$ and is, to leading order in $s$, advected by $\uE$ alone in the same way. This is consistent with \citet{Schekochihin2019}, who identified a passively advected kinetic field at low $\beta$ with the kinetic analogue of the entropy mode.
